\chapter{来源与史料说明}
\label{app:sources}

本附录不是把材料排成固定的高低次序，而是说明不同材料最适合回答什么、又不能回答什么。本卷事件账本的每条记录均保留核心来源、现代研究方向和可信度说明；正文只从已入账事件发展叙事。

\section{正史与编年：提供骨架，不替代现场}

\begin{description}
  \phantomsection\label{src:jinshu}\sourceindexentry{晋书}{西晋、东晋及十六国的主干材料，帝纪、志、列传和载记各有不同用途。唐初编纂、人物褒贬、谶纬与以晋为中心的命名，要求与其他材料交读。}
  \phantomsection\label{src:songshu}\sourceindexentry{宋书}{刘宋政治、北伐、州郡、礼制和财政的基础入口；江南朝廷视角不能代替北方与地方社会。}
  \phantomsection\label{src:nanqishu}\sourceindexentry{南齐书}{短朝继承、军政与萧氏叙事的重要入口；作者与皇室的关系尤其需要辨认。}
  \phantomsection\label{src:liangshu}\sourceindexentry{梁书}{梁武帝、侯景之乱、使节与南方崩解的主干；亡国叙事容易把州镇、财政和道路问题人格化。}
  \phantomsection\label{src:chenshu}\sourceindexentry{陈书}{陈的建国、江南重建与隋前关系的主要材料；须以北朝史、墓志和考古补足其窄小的朝廷视野。}
  \phantomsection\label{src:weishu}\sourceindexentry{魏书}{北魏建国、迁都、均田、三长、官制与释老史料的主干；“秽史”争论以及族名、政策与执行之间的距离均需正视。}
  \phantomsection\label{src:beiqishu-zhoushu}\sourceindexentry{北齐书与周书}{东魏北齐、关中西魏北周与灭齐路径的基础；唐代回望可能加强北周—隋连续性。}
  \phantomsection\label{src:suishu}\sourceindexentry{隋书}{五八九年统一以及制度志的后出校核入口；不得用后来的定制覆盖早期异文。}
  \phantomsection\label{src:zizhi-tongjian}\sourceindexentry{资治通鉴}{跨政权年月日和前后关联的编年参照；它是北宋的选择与重编，不能替代所据正史、文书或物质材料。}
\end{description}

\section{文书、碑刻、宗教与遗址}

\phantomsection\label{src:material-record}
\begin{description}
  \item[墓志、墓券与碑刻] 可校正少数人物的日期、官职、籍贯、婚姻和自我表述。追赠、谱系和道德修辞不应被直接外推为社会统计。
  \item[造像记、石窟和寺院遗址] 可见供养、工匠、材料和宗教活动的局部网络。题记日期、洞窟营造期和文本成书期常不相同；一处遗址不代表全体信众。
  \item[佛道文献] 僧传、经录、护教辩难、道经和旅行记录有助于理解社群怎样组织知识和记忆。它们的典范化与神异叙事不能充当政策执行或人口的统计。
  \item[河西、西域及城市材料] 吐鲁番、高昌、楼兰/尼雅文书、城址、军镇、道路、仓储和墓葬，能检验地方行政与交通。语言、纪年、出土环境和政权归属均须逐案核对。
\end{description}

\section{怎样按问题回到材料}
\phantomsection\label{src:question-method}

战争与外交先看本纪、列传和编年，再以河流、道路、城址及对方材料限制兵数和路线；制度先看志书、诏令与官名，再追问谁在地方登记、征发和执行；迁徙先区分侨置名目、墓志自述和可见的居住网络；宗教先分开国家政策、教团文本、个人供养和物质遗存。读者可从\hyperref[app:index]{“事件与图件索引”}回到正文，再参看\hyperref[app:controversies]{“争议怎样改变读法”}。

\section{基础书目与读法索引}

以下条目只列本册已用于建立年序、制度问题或材料边界的书目入口。现有材料记录能够核对书名、作者/编者与用途，但没有为全卷统一登记所据点校本、出版年或页码；因此不把任何一条写成精确引文的替代物。需要核页或核措辞时，应回到事件所用篇目、可靠底本和出土材料的释读报告。

\begin{description}
  \item[《晋书》，房玄龄等撰] 适用于西晋、东晋与十六国的年序、官名和叙事分歧；本册以它建立问题和交叉路径，不能单独裁定人物动机、兵数或十六国疆界。
  \item[《宋书》，沈约撰；《南齐书》，萧子显撰] 适用于南朝继承、州镇与北伐的叙事骨架；它们不能代替地方墓志、城市遗址和北方材料。
  \item[《梁书》《陈书》，姚思廉撰] 适用于梁陈政局、侯景之乱与江南重建；亡国叙事不应覆盖州镇、财政和道路问题。
  \item[《魏书》，魏收撰；《北齐书》，李百药撰] 适用于北朝军政、制度、迁都与宗教史料的正史入口；“汉化”和正统仍须分开文本框架与地方执行。
  \item[《周书》，令狐德棻等撰；《隋书》，魏徵等撰] 适用于北周、隋初与再统一；不能以唐初的回望直接推论北周—隋连续性的全部过程。
  \item[《资治通鉴》，司马光编] 适用于跨政权的年月日、前后顺序和事件链追踪；它是北宋重编的编年叙事，不能代替所据正史、文书、题记或物质遗存。
  \item[《建康实录》，许嵩撰；《水经注》，郦道元撰] 分别为建康及南方政治地理、河流道路与区域空间提供补充入口；两书的成书时间、传本和所引旧闻须逐项核验，不能直接复原战场路线或地方社会全貌。
  \item[墓志、造像记、吐鲁番与河西文书、城址报告] 适用于日期、官职、供养、运输和地方行政的局部检验；其出土环境、保存偏差与个案性质禁止把少数材料外推为全境统计。
\end{description}
